%%%%%%%%%%%%%%%%%%%%%%%%%%%%%%%
\chapter{Redesign Incompressible Navier--Stokes Solver}
\label{chap:incns}

% Local notation for this chapter. The names are prefixed with "incns" so
% that they do not clash with the developer guide preamble.
\newcommand{\incnsv}[1]{\boldsymbol{#1}}
\newcommand{\incnsdt}{\Delta t}
\newcommand{\incnsCC}{\nabla\times\nabla\times}

% Long identifiers in typewriter font cannot be hyphenated; allow looser
% spacing in this chapter instead of overfull lines. Reset at the end.
\sloppy

This chapter summarises \code{IncNavierStokesSolverRedesign}, the
incompressible Navier--Stokes solver built on the redesign operators of
the \code{MultiRegions} and \code{SolverCore} libraries. It covers the time
discretisation, the boundary conditions specific to the solver, initial
conditions and restarts, a known start-up pressure transient, and the
session parameters. Two further chapters and a standalone implementation
note in \path{docs/developer-guide/redesign/} give the full detail; the note
is referred to where useful:
\begin{itemize}
  \item Chapter~\ref{chap:hopressure-bndcond}: the higher-order pressure
  boundary condition (HOPBC);
  \item Chapter~\ref{chap:houtflow-bndcond}: the Dong outflow condition
  (HOutflow).
\end{itemize}

%%%%%%%%%%%%%%%%%%%%%%%%%%%%%%%
\section{Overview}
\label{sec:incns-overview}

The solver implements the velocity-correction scheme (VCS) of Karniadakis,
Israeli and Orszag~\cite{Karniadakis1991} for continuous Galerkin
(\code{Projection=Galerkin}) discretisations. Each time step solves a
pressure Poisson problem followed by one Helmholtz (semi-implicit) or linear
advection--diffusion--reaction (ADR, linear-implicit) problem per velocity
component. All element, boundary and linear-system work is done by
redesign operators, so the same solver code runs on every execution space
(\code{Serial}, \code{AVX}, \code{Device} for CUDA or HIP) and with every
element-operator implementation (\code{StdMat}, \code{SumFac}, and
\code{SumFacTOP} on \code{Device} only). There are no backend-specific
branches in the solver.

The main pieces are listed below; paths are relative to
\path{solvers/IncNavierStokesSolverRedesign/} or \path{library/SolverCore/}.

\begin{itemize}
  \item \path{EquationSystems/VelocityCorrectionScheme.{h,cpp}}: the
  equation system, registered as \code{VelocityCorrectionScheme}. It derives
  from \code{SolverCore::UnsteadySystem} and provides the two time-integrator
  callbacks \code{EvaluateAdvection\_SetPressureBCs} (explicit part) and
  \code{SolveUnsteadyStokesSystem} (implicit part).
  \item \path{BoundaryConditions/HOPressureBndCond/}: the
  \code{HOPressureBndCondOp}, a Neumann boundary-condition operator added to
  the pressure Poisson solve (Section~\ref{sec:incns-hopbc}).
  \item \path{BoundaryConditions/HOutflowBndCond/}: the
  \code{HOutflowBndCondOp}, which holds the boundary data shared by a pressure
  Dirichlet and a velocity Neumann condition and presents them to the linear
  systems through two adapters (Section~\ref{sec:incns-houtflow}).
  \item \path{TimeOps/IMEX} and \path{TimeOps/IMEXdirk} in \code{SolverCore}:
  the IMEX multistep scheme and the IMEX DIRK scheme used for its start-up.
  \item \path{GlobalLinSysOps/LinearSystems} in \code{SolverCore}:
  \code{PoissonSolveOp} (pressure), \code{HelmSolveOp} (semi-implicit
  velocity), \code{LinearADRSolveOp} (linear-implicit velocity) and
  \code{FwdTransOp} (projections); \path{GlobalLinSysOps/LinearSolvers}: the
  iterative solvers (\code{ConjGrad}, \code{GMRES} and others);
  \path{PreconOps}: the \code{Diagonal} and \code{Null} preconditioners.
  \item \path{Filters} in \code{SolverCore}: the filter available to the
  solver is \code{Checkpoint}. Checkpoints are written only by the
  \code{Checkpoint} filter; \code{IO\_CheckSteps} and \code{IO\_CheckTime}
  are ignored with a warning.
\end{itemize}

The variables must be ordered as velocity components, optional passive
scalars, and pressure last. The pressure has its own \code{ContField}, which
shares the velocity expansion when the session gives both the same
expansion. The pressure is excluded from time integration: the time operator
advances the velocity (and scalar) fields only.

%%%%%%%%%%%%%%%%%%%%%%%%%%%%%%%
\section{Governing equations and time discretisation}
\label{sec:incns-time}

The solver integrates
\begin{equation}
  \frac{\partial\incnsv u}{\partial t} = \incnsv N(\incnsv u) - \nabla p
  + \nu\nabla^2\incnsv u,\qquad \nabla\cdot\incnsv u = 0,\qquad
  \incnsv N(\incnsv u) = -(\incnsv u\cdot\nabla)\incnsv u + \incnsv f,
  \label{eq:incns-ns}
\end{equation}
with kinematic viscosity $\nu$ (\code{Kinvis}) and body forcing $\incnsv f$
from the session \code{FORCING} block. Only the convective form of the
advection term is implemented. With \code{SPECTRALHPDEALIASING} (on by
default) it is evaluated with the 3/2-rule \code{AdvectionDealiasOp},
otherwise with \code{AdvectionOp}.

\subsection{IMEX scheme and the implicit-stage size}

The IMEX scheme of order $J$ (\code{METHOD=IMEX}, \code{ORDER=}$J$) is the
stiffly stable scheme of~\cite{Karniadakis1991}. The implementation divides
it by the coefficient $\gamma_0$ of $\incnsv u^{n+1}$, and the implicit
callback receives
\begin{equation}
  \hat{\incnsv u} = \sum_{q=0}^{J-1}\frac{\alpha_q}{\gamma_0}\incnsv u^{n-q}
  + \incnsdt\sum_{q=0}^{J-1}\frac{\beta_q}{\gamma_0}\incnsv N^{n-q},
  \qquad \lambda\incnsdt = \frac{\incnsdt}{\gamma_0},
  \label{eq:incns-uhat}
\end{equation}
where $\lambda=1/\gamma_0$ is \code{m\_gamma} in \code{IMEXOpImpl}. The
fourth argument of \code{SolveUnsteadyStokesSystem}, \code{dt\_inv\_gamma}, is
the implicit-stage size $\lambda\incnsdt$. The explicit
callback stores $\incnsdt\,\incnsv N$: the advection operator is applied
with scale $-\incnsdt$ and the forcing with scale $\incnsdt$.

\begin{table}[htbp]
\centering
\small
\begin{tabular}{@{}lcccc@{}}
\toprule
order $J$ & 1 & 2 & 3 & 4 \\
\midrule
$\lambda=1/\gamma_0$ & $1$ & $2/3$ & $6/11$ & $12/25$ \\
weights $c_q=\beta_q/\gamma_0$ & $[1]$ & $[4,\,-2]/3$ &
$[18,-18,6]/11$ & $[48,-72,48,-12]/25$ \\
start-up scheme & none & IMEXdirk 2\_3\_2 & IMEXdirk 3\_4\_3 &
IMEXdirk 3\_4\_3 \\
\bottomrule
\end{tabular}
\caption{IMEX coefficients (\code{GetIMEXCoefficients} in
\code{IMEXKernelLaunchers.hpp}) and start-up schemes. In every multistep step
$\sum_q c_q=\lambda$.}
\label{tab:incns-imex}
\end{table}

\textbf{Start-up.} A multistep scheme of order $J>1$ has no history at
the first step. While fewer than $J$ steps have been taken,
\code{IMEXOpImpl::v\_Apply} advances one step with an \code{IMEXdirk} scheme
of order $\min(3,J)$ (Table~\ref{tab:incns-imex}), to which it copies its
callbacks. For $J=2$ this is IMEXdirk 2\_3\_2 with two implicit stages, both
of size $\lambda_s\incnsdt$ with $\lambda_s=1-1/\sqrt2\approx0.2929$. The
history is not written to checkpoints, so the start-up step is repeated after
every restart (Section~\ref{sec:incns-kick}).

\textbf{Explicit-contribution extrapolation.} The VCS needs
extrapolations of quantities other than the solution, combined with the same
weights as the explicit term of the active step or stage. The solver calls
\code{EnableExplicitContributionExtrapolation()} on the time operator. Before
each implicit solve, the operator publishes the active weights $c_i$: the
first $J$ IMEX coefficients in a multistep step, or the explicit stage
coefficients of the current DIRK stage. The solver keeps up to three history
deques per history id, as deep as the implicit solves of the scheme combine
explicit contributions (\code{GetExplicitContributionHistorySize}), refreshed
in \code{EvaluateAdvection\_SetPressureBCs}:
\begin{itemize}
  \item the velocity history, $\incnsv u$, for the advection velocity of the
  linear-implicit formulation and, on a rank with a higher-order pressure or
  outflow region, for the boundary data;
  \item the pressure-BC RHS history, the explicit forcing $\incnsv N$, only on
  a rank with a higher-order pressure region whose velocity is not Dirichlet
  in every component (Section~\ref{sec:incns-hopbc});
  \item the advection history, $-\incnsdt(\incnsv u\cdot\nabla)\incnsv u$
  (linear-implicit formulation only).
\end{itemize}
Every start-up scheme has a history id of its own; its histories are freed
when the main scheme evaluates its explicit term again
(\code{FreeStartupHistories}).
\code{TimeOp::ExtrapolateExplicitContribution} forms
$\incnsv X^*=\sum_i c_i\,\incnsv h_i$ from a history $\incnsv h$. The
advection history is combined with the weights $c_i$ themselves, so that it
matches the explicit contribution held in $\hat{\incnsv u}$. For the
boundary data, the velocity and pressure-BC RHS histories are combined with the weights scaled to sum to $\lambda\incnsdt$,
the scale being $\lambda\incnsdt/\sum_ic_i$ with $\sum_ic_i$ from
\code{TimeOp::GetExplicitContributionCoefficientSum}, so that with
$s=1/(\lambda\incnsdt)$ the velocity history gives
\begin{equation}
  \incnsv u^* = s\,\incnsv V,\qquad
  \incnsv V=\frac{\lambda\incnsdt}{\sum_i c_i}\sum_i c_i\,\incnsv u_i,
  \qquad
  \incnsv u^* = \sum_i\frac{c_i}{\sum_j c_j}\incnsv u_i ,
  \label{eq:incns-ustar}
\end{equation}
an extrapolation with weights summing to one. In every multistep step
$\sum_i c_i=\lambda$ (e.g.\ $\incnsv u^*=2\incnsv u^n-\incnsv u^{n-1}$ for
$J=2$), as in DIRK stage 1, and the scale is $\incnsdt$; in DIRK stage 2 it
is not (Section~\ref{sec:incns-kick}).

\subsection{The velocity-correction stage}
\label{sec:incns-vcs}

Given $\hat{\incnsv u}$ and $\lambda\incnsdt$,
\code{SolveUnsteadyStokesSystem} performs the following steps.

\begin{enumerate}
  \item \textbf{Pressure forcing.} $f_p=\nabla\cdot\hat{\incnsv u}/(\lambda
  \incnsdt)$ with \code{DivergenceOp}. The explicit advection stays in
  $\hat{\incnsv u}$ for both formulations.
  \item \textbf{Pressure boundary data.} Extrapolate the pressure-BC
  histories, set the kinematic viscosity, the time $t^{n+1}$ and the scale
  $s=1/(\lambda\incnsdt)$ on the boundary operators (\code{SetKinvis},
  \code{SetTime}, \code{SetScale}), and update the Neumann data of
  \code{HOPressureBndCondOp} (\code{UpdateWithAcceleration},
  Section~\ref{sec:incns-hopbc}) and the Dirichlet data of
  \code{HOutflowBndCondOp} (\code{UpdatePressure},
  Section~\ref{sec:incns-houtflow}).
  \item \textbf{Singular pressure system.} If every pressure boundary is
  Neumann or periodic on all ranks (\code{v\_GetSystemSingularChecks}),
  \code{MeanRemovalOp} removes the mean of $f_p$, which makes a periodic
  system compatible. A Neumann boundary with non-zero data, such as a
  higher-order pressure region, adds a surface integral that the mean
  removal does not balance.
  \item \textbf{Pressure solve.} Solve the strong pressure Poisson problem
  \begin{equation}
    \nabla^2p^{n+1}=\frac{1}{\lambda\incnsdt}\nabla\cdot\hat{\incnsv u}
    \ \text{in }\Omega,\qquad
    \frac{\partial p^{n+1}}{\partial n}=g\ \text{on }\Gamma_H,\qquad
    p^{n+1}=p_D\ \text{on }\Gamma_D,
    \label{eq:incns-poisson}
  \end{equation}
  with \code{PoissonSolveOp} and a \code{ConjGrad} solver, then
  \code{BwdTransOp} and \code{PhysDerivOp} for $p^{n+1}$ and $\nabla
  p^{n+1}$.
  \item \textbf{Velocity solve.} For the semi-implicit formulation solve, per
  component,
  \begin{equation}
    \nabla^2\incnsv u^{n+1}-\frac{1}{\nu\lambda\incnsdt}\incnsv u^{n+1}
    = \frac{1}{\nu}\Bigl(\nabla p^{n+1}
      - \frac{\hat{\incnsv u}}{\lambda\incnsdt}\Bigr)
    \label{eq:incns-helm}
  \end{equation}
  with \code{HelmSolveOp} and \code{ConjGrad}. For the linear-implicit
  formulation the extrapolated advection is removed from the right-hand
  side, $\hat{\incnsv u}_L=\hat{\incnsv u}-\sum_ic_i\incnsv h_i^{\mathrm{adv}}$,
  and the linearised advection moves to the operator:
  \begin{equation}
    \nabla^2\incnsv u^{n+1}
    - \frac{1}{\nu}(\incnsv U\cdot\nabla)\incnsv u^{n+1}
    - \frac{1}{\nu\lambda\incnsdt}\incnsv u^{n+1}
    = \frac{1}{\nu}\Bigl(\nabla p^{n+1}
      - \frac{\hat{\incnsv u}_L}{\lambda\incnsdt}\Bigr),
    \label{eq:incns-adr}
  \end{equation}
  solved with \code{LinearADRSolveOp} and \code{GMRES}. The right-hand side
  lost the explicit advection $\incnsdt\sum_ic_i(\incnsv u_i\cdot\nabla)
  \incnsv u_i$, which the operator puts back linearised about the
  extrapolated velocity $\incnsv u^*$ of \eqref{eq:incns-ustar}, as
  $\incnsdt\sum_ic_i\,(\incnsv u^*\cdot\nabla)\incnsv u^{n+1}$. It weights
  the advection by $\lambda\incnsdt$, so its advection velocity is
  \begin{equation}
    \incnsv U=\frac{\sum_ic_i}{\lambda}\,\incnsv u^*
    =\frac{1}{\lambda}\sum_ic_i\,\incnsv u_i ,
    \label{eq:incns-advvel}
  \end{equation}
  formed from the velocity history by
  \code{TimeOp::ExtrapolateExplicitContribution} with the scale
  $\incnsdt/(\nu\lambda\incnsdt)$ (\code{m\_advVel} holds
  $\incnsv U/\nu$); the sum of the weights cancels. In a multistep step
  $\sum_ic_i=\lambda$ and $\incnsv U=\incnsv u^*$ extrapolates the velocity
  to $t^{n+1}$ at the order $J$ of the scheme
  ($2\incnsv u^n-\incnsv u^{n-1}$ for $J=2$). In a DIRK stage the weights
  have another sum (Section~\ref{sec:incns-kick}) and $\incnsv u^*$ is their
  weighted mean of the earlier stage values, accurate to first order only:
  every start-up step has a second-order local error, which limits a
  multistep scheme to second order (Section~\ref{sec:incns-weak-li}), and an
  IMEXdirk scheme run on its own advects to first order. Before the solve, \code{SetPressure} hands
  the solved pressure to \code{HOutflowBndCondOp}, whose velocity Neumann
  adapter forms the outflow flux from it when the solve applies it.
\end{enumerate}

The preconditioners of the pressure and semi-implicit velocity systems are
cached per value of $\lambda\incnsdt$, so the start-up stages and the
multistep steps each build theirs once. The linear-implicit velocity
preconditioner is rebuilt every stage, because $\incnsv U$ changes.

\subsection{Weak pressure form}
\label{sec:incns-weak}

The weak form of the pressure equation is selected with
\code{SOLVERINFO PressureFormulation=Weak} (case insensitive), independently
of \code{FORMULATION}. The default is \code{Strong} with semi-implicit and
\code{Weak} with linear-implicit time-stepping. The redesign
Poisson operator solves $(\nabla\phi,\nabla p)_\Omega=F+\int_{\Gamma_N}\phi
g\,\mathrm dS$ for a forcing $F$ and a Neumann datum $g$. The strong form takes
$F=-(\phi,\nabla\cdot\incnsv f_w)_\Omega$ with
$\incnsv f_w=\hat{\incnsv u}/(\lambda\incnsdt)$; integrating it by parts,
\begin{equation*}
  -(\phi,\nabla\cdot\incnsv f_w)_\Omega=(\nabla\phi,\incnsv f_w)_\Omega
  -\int_{\partial\Omega}\phi\,\incnsv n\cdot\incnsv f_w\,\mathrm dS ,
\end{equation*}
and moving the boundary term into the Neumann datum of the HOPBC
(\eqref{eq:incns-hopbc-discrete}, whose $s\incnsv P$ is $\incnsv f_w$) gives
\begin{equation}
  (\nabla\phi,\nabla p^{n+1})_\Omega = (\nabla\phi,\incnsv f_w)_\Omega
  - \int_{\Gamma_H}\phi\,\incnsv n\cdot\Bigl(
    \frac{\incnsv u_D^{n+1}}{\lambda\incnsdt}
    + \nu\,\incnsCC\incnsv u^*\Bigr)\mathrm dS .
  \label{eq:incns-weak}
\end{equation}
The predictor, explicit advection and forcing enter only through the volume
term $(\nabla\phi,\incnsv f_w)$, computed with
\code{MultiRegions::IProductWRTDerivBaseOp<FieldState::Coeff>} and passed to
\code{PoissonSolveOp::ApplyRhs}, which takes the right-hand side in
coefficient form, adds the Neumann data and the Dirichlet lift to it and
solves; the physical-space \code{Apply} forms its right-hand side and then
solves the same way. Higher-order pressure regions then need only the viscous
term and the boundary acceleration, which
\code{HOPressureBndCondOp::UpdateWeak} provides: the flux of
\code{UpdateWithAcceleration} less $s\,\incnsv n\cdot\incnsv P$, that is
$g=-s\,\incnsv n\cdot(\incnsv u_D+\nu\incnsCC\incnsv V)$. A velocity
component that is not Dirichlet on a region takes $\incnsv P-\incnsv R$, the
velocity the predictor carries, as its boundary velocity, as in the strong
form, so its flux there is $-s\,n_d(\nu C_d(\incnsv V)+P_d-R_d)$; the
explicit-forcing history is kept for such regions only, as in the strong
form. Further choices:
\begin{itemize}
  \item semi-implicit: $\incnsv u^*$ is extrapolated as in the strong form;
  \item linear-implicit: the advection in $\incnsv f_w$ and the velocity in
  the curl--curl and HOutflow terms are taken at the latest level, weighted by
  $\sum_ic_i$ (\code{TimeOp::LatestExplicitContribution}), for stability at
  large time steps; this lowers the time order of the pressure near
  higher-order pressure regions (Section~\ref{sec:incns-weak-li});
  \item the HOutflow pressure is a Dirichlet value, on which the test
  functions vanish, and is the same in both forms;
  \item a Neumann pressure region without the higher-order tag keeps its
  session value $g$ as the datum of the weak form, which therefore prescribes
  $(\nabla p-\incnsv f_w)\cdot\incnsv n=g$ there, not
  $\partial p/\partial n=g$;
  \item the volume term is orthogonal to the constants, so a singular
  pressure system without higher-order pressure regions, such as a periodic
  one, is compatible as it stands and no mean is removed. With higher-order
  pressure regions the Neumann flux has a non-zero surface integral, which
  the weak form has no correction to balance yet, so this case is refused at
  set-up;
  \item velocity and pressure must share their elements, blocks and
  quadrature points, checked at set-up with
  \code{ExpList::SameQuadratureLayout},
  since the volume term reads the velocity right-hand side on the pressure
  expansion.
\end{itemize}

\subsubsection{Linear-implicit formulation: stability and time order}
\label{sec:incns-weak-li}

The linear-implicit velocity solve advects with $\incnsv U$ of
\eqref{eq:incns-advvel}, so its time step may exceed the advective CFL
limit. In a multistep step $\incnsv U=\incnsv u^*$ differs from
$\incnsv u^{n+1}$ by $O(\incnsdt^J)$; in the IMEXdirk start-up stages it is
a mean of stage values, which differs by $O(\incnsdt)$ and limits the
formulation to second order in time for $J=3$ on a flow whose advection
changes in time. The pressure step still treats the advection in
$\incnsv f_w$, the curl--curl term and the HOutflow velocity explicitly.
Extrapolating them, as the strong and the semi-implicit weak forms do,
brings back a time-step limit: the extrapolated
$\incnsv u^*=\sum_ic_i\incnsv u_i/\sum_ic_i$ amplifies the high-frequency
part of the history (by 3 for $J=2$ and 7 for $J=3$). The weak form takes
them at the latest level instead; the price is an $O(\incnsdt)$ error in
the explicit terms of its pressure step, which makes the form second order
in time for $J=3$ on a flow whose advection does not change in time, and
leaves its pressure near higher-order pressure regions close to first order
on a flow whose advection does (below).

On the manufactured flow \path{HOPressureTimeConvergence2D.xml} (8, 16 and
32 steps, Serial and AVX SumFac; its advection vanishes, so only the
boundary terms matter) the final rates of
$u$, $v$, $p$ for $J=3$ are $1.98$, $1.98$, $1.60$ with the latest level and
$3.22$, $3.23$, $2.92$ with the curl--curl velocity extrapolated; for $J=2$
they are $2.06$, $2.22$, $1.84$ and $2.17$, $2.22$, $1.94$. The level of the
advection in $\incnsv f_w$ does not change these rates. The strong form is
third order for $J=3$ ($3.23$, $3.25$, $2.93$). The largest stable
time step (400 steps from a perturbed or impulsive start, Serial SumFac) is,
with the values of the earlier advection velocity $\incnsv U=\incnsv u^n$ in
parentheses where they differ:
\begin{center}\small
\begin{tabular}{@{}lcccc@{}}
\toprule
case, $J$ & strong & weak & weak, curl--curl $\incnsv u^*$\textsuperscript{a} &
weak, all extrapolated\textsuperscript{a} \\
\midrule
\path{KovasznayFlowHOPBC2D}, 2 & $0.01$ & $0.3$ ($0.35$) & $0.35$ & $0.01$ \\
\path{KovasznayFlowHOPBC2D}, 3 & $0.005$ & $0.2$ ($0.25$) & $0.1$ & $<0.005$ \\
\path{CylinderFlowHOPBC2D}, 2 & $0.6$\textsuperscript{b} & $\ge3$ & $\ge3$ & $0.5$ \\
\path{CylinderFlowHOPBC2D}, 3 & $0.25$\textsuperscript{b} & $\ge3$ & $0.75$ & $0.2$ \\
\path{HOutflowKovasznay2D}, 2 & $0.01$ & $0.8$ ($1$) & $0.02$ & $0.01$ \\
\path{HOutflowKovasznay2D}, 3 & $<0.005$ & $0.35$ ($0.4$) & & \\
\path{HOutflowBackflow2D}, 2 & $0.01$ & $0.5$ ($1$) & $0.01$ & $0.01$ \\
\path{HOutflowBackflow2D}, 3 & $<0.005$ & $0.4$ ($0.8$) & & \\
\bottomrule
\end{tabular}

\textsuperscript{a} measured with $\incnsv U=\incnsv u^n$;
\textsuperscript{b} $0.5$ and $0.2$ on the coarser grid of time steps used
for $\incnsv U=\incnsv u^n$.
\end{center}
In the HOutflow cases the last column also extrapolates the HOutflow
velocity, which shares the history. Extrapolating the curl--curl velocity
only would restore third order on the flow above but reduce the stable time
step by a factor of 2.5 to more than 4 for $J=3$. For $J=2$ it costs no
stability in the higher-order pressure cases but gains little, so the latest
level is kept for every order. Advecting with $\incnsv u^*$ instead of
$\incnsv u^n$ costs little stability with higher-order pressure regions and
up to a factor of two with HOutflow regions. In \path{HOutflowKovasznay2D}
the weak form with $\incnsv u^*$ also keeps the error of the steady solution
only up to a time step of $0.12$ ($J=2$) and $0.15$ ($J=3$), and stays
bounded but inaccurate above it, where with $\incnsv u^n$ it kept that error
up to its stability limit. \code{VelocityCorrectionScheme::InitialiseParameters()}
warns (on the root rank) when \code{FORMULATION=LinearImplicit} is used with
a time integration scheme of order three or higher, whose order it cannot
reach: the weak form is the more stable, with a lower-order pressure near
higher-order pressure regions, and \code{PressureFormulation=Strong} the more
accurate but much less stable.

\textbf{Advective time order.} The flow above has no advection, so it does
not see the advection velocity $\incnsv U$ of the velocity solve.
\path{check_hopbc_time_order.py --manufactured advective} sets the
unsteady flow
$\incnsv u=\cos t\,(x^2,-2xy)+\sin t\,(2xy,-y^2)$ on the same session, with
$(\partial_t\incnsv u\cdot\nabla)\incnsv u$ non-zero and not a gradient; it
is quadratic, so the spatial error is zero and only the time error remains.
The final rates of $u$, $v$, $p$ (8, 16 and 32 steps to $t=0.2$, Serial and
AVX SumFac, \code{NUMMODES} 8) are:
\begin{center}\small
\begin{tabular}{@{}lcc@{}}
\toprule
formulation, pressure form & $J=2$ & $J=3$ \\
\midrule
semi-implicit, strong & $2.22$, $2.21$, $1.93$ & $2.89$, $2.79$, $2.95$ \\
semi-implicit, weak & $2.22$, $2.21$, $1.93$ & $2.89$, $2.78$, $2.98$ \\
linear-implicit, strong & $2.26$, $2.32$, $1.92$ & $1.98$, $2.18$, $2.02$ \\
linear-implicit, weak & $2.03$, $1.99$, $1.63$ & $1.98$, $1.95$, $1.48$ \\
\bottomrule
\end{tabular}
\end{center}
With $\incnsv U=\incnsv u^n$ every linear-implicit rate was between $0.95$
and $1.04$. The linear-implicit formulation is second order for $J=3$ as
well, because of the start-up stages. The weak-form pressure falls towards
first order under refinement: from 32 to 64 and from 64 to 128 steps its
rates are $1.51$ and $1.37$ for $J=2$, $1.36$ and $1.26$ for $J=3$, while
the strong form stays at $1.94$ to $2.02$. Extrapolating the advection in
$\incnsv f_w$ as well raises the weak-form rates from 32 to 64 steps
(\code{NUMMODES} 6) from $1.30$ and $1.19$ to $1.80$ and $1.31$, and
extrapolating also the curl--curl velocity gives the rates of the strong
form, at the cost of the stable time step above. The ctests
\code{IncNavierStokesSolverRedesign\_HOPBC\_Advective\_TimeIntegrationOrders23}
(semi-implicit, both forms, at the order of the scheme),
\code{\_Advective\_LinearImplicit\_TimeIntegrationOrders23} (strong form, every
rate at least $1.7$ for $J=2$ and $3$) and
\code{\_Advective\_WeakPressure\_LinearImplicit\_TimeIntegrationOrders23}
(velocity rates at least $1.7$, pressure at least $1.1$, which rules out
first order) check these rates. An IMEXdirk scheme run on its own advects
with stage means in every step, and the linear-implicit velocity is first
order on this flow with it (\code{IMEXdirk23} and \code{IMEXdirk34}; the
pressure of every formulation is first order with these schemes too).
The regression cases are steady or use first-order IMEX, for which
$\incnsv U=\incnsv u^n$. The manufactured flow of
\path{check_time_order.py},
$\incnsv u=(1+0.1\sin t)\,\incnsv u_0(\incnsv x)$, changes its advection by
$0.1\cos t\,(\incnsv u_0\cdot\nabla)\incnsv u_0$ only, at $\nu=1$ and with 2
and 4 steps, where the error does not show, and
\path{check_imexdirk_forced_mms_2d.py} checks that four IMEXdirk steps of
the same kind of flow run, not a rate.

Legacy \code{VCSImplicit} advects with the extrapolated velocity
$\incnsv u^*$ by default (\code{SOLVERINFO AdvectionVelocity=Extrapolated};
\code{Updated} takes a first-order predictor of $\incnsv u^{n+1}$ instead),
as the redesign does, and takes the advection of its weak pressure forcing
and the curl--curl velocity of its higher-order pressure flux at the latest
level, as the redesign's linear-implicit weak form does. On this session it
is first order in the velocity, with pressure rates of $0.4$ to $0.7$, for
both flows, including the one without advection, so from its higher-order
pressure condition; with a Dirichlet pressure on every edge it is second
order for $J=2$ and $J=3$ in both of its advection velocity choices (its
semi-implicit scheme is also second order only for $J=3$ there). Its stable
time steps (legacy \code{VCSImplicit}, \code{Extrapolated}) are $\ge2$ in
\path{KovasznayFlowHOPBC2D} and \path{CylinderFlowHOPBC2D}, $0.02$ and
$0.005$ in \path{HOutflowKovasznay2D} for $J=2$ and $3$, and $0.01$ and
below $0.005$ in \path{HOutflowBackflow2D}.

%%%%%%%%%%%%%%%%%%%%%%%%%%%%%%%
\section{Boundary conditions}
\label{sec:incns-bc}

Ordinary Dirichlet, Neumann and Robin conditions on velocity and pressure
are imposed by the generic \code{DirBndCondOp}, \code{NeuBndCondOp} and
\code{RobBndCondOp} of the linear-system operators: Dirichlet values are set
in the coefficient solution, Neumann data is added to the right-hand side as
a boundary inner product, and Robin terms enter the operator of the
iterative solver. The two solver-specific conditions below are selected with
\code{USERDEFINEDTYPE} and plug into the same interfaces:
\code{m\_poissonSolveOp->AddNeumannBCOp},
\code{m\_poissonSolveOp->AddDirichletBCOp} and
\code{m\_fieldsSolveOp->AddNeumannBCOp} in
\code{VelocityCorrectionScheme::v\_InitialiseOperators}. Each operator is
built only for a session with a region of its type
(\code{HasUserDefinedType}, reduced over the ranks). A Neumann operator is
added with \code{LinearSystemOp::AddNeumannBCOp(op, userDefinedType)}, which
makes the default Neumann operator of the system leave the regions of that
type out (\code{NeuBndCondOp::ExcludeUserDefinedType}).

\subsection{Higher-order pressure condition (HOPBC)}
\label{sec:incns-hopbc}

\textbf{Continuous condition.} Writing the viscous term in rotational
form, $\nabla^2\incnsv u=-\incnsCC\incnsv u$ for $\nabla\cdot\incnsv u=0$, and
taking the normal component of~\eqref{eq:incns-ns} gives
\begin{equation}
  \frac{\partial p}{\partial n}=\incnsv n\cdot\Bigl(\incnsv N
  - \nu\,\incnsCC\incnsv u - \frac{\partial\incnsv u}{\partial t}\Bigr).
  \label{eq:incns-hopbc}
\end{equation}

\textbf{Time-discrete form.}
\code{HOPressureBndCondOp::UpdateWithAcceleration} receives the extrapolated
explicit forcing $\incnsv R$, the extrapolated velocity $\incnsv V$ (both
scaled as in~\eqref{eq:incns-ustar}) and the predictor
$\incnsv P=\hat{\incnsv u}$; the scale $s=1/(\lambda\incnsdt)$, the
viscosity and the time are set beforehand with \code{SetScale},
\code{SetKinvis} and \code{SetTime}. The explicit parts of the base flux
$s(\incnsv R-\nu\incnsv C)$ and the acceleration flux
$s(\incnsv P-\incnsv R)$ cancel, so the operator evaluates
\begin{equation}
  g = s\,\incnsv n\cdot\bigl[\incnsv P-\incnsv u_D(t^{n+1})
      -\nu\,\incnsv C(\incnsv V)\bigr],\qquad
  \incnsv C(\incnsv V)=\incnsCC\incnsv V,
  \label{eq:incns-hopbc-discrete}
\end{equation}
where $\incnsv u_D$ is the discrete Dirichlet velocity: the session
expressions evaluated on the boundary (\code{ExpressionOp}) and projected
onto the boundary expansion of the velocity (\code{FwdTransBCOp}, then
\code{BwdTransOp}), as the velocity solve imposes it. Its coefficients are
kept and projected once, unless a Dirichlet velocity condition of the region
is time-dependent; a point region takes the expression itself. A set-up \code{ASSERTL0} checks
that the velocity and pressure boundary expansions share their points.
$s(\incnsv P-\incnsv u_D)$ is the negative discrete acceleration plus the
explicit term, and vanishes at a steady state; $s\incnsv V=\incnsv u^*$
of~\eqref{eq:incns-ustar}. A velocity component $d$ that is not Dirichlet on
a region has no prescribed velocity and takes no acceleration: its flux
there is $s\,n_d[R_d-\nu C_d(\incnsv V)]$, which is why $\incnsv R$ is kept
for such regions only. \code{Update}, without the predictor, evaluates
$g=s\,\incnsv n\cdot[\incnsv R-\nu\incnsv C(\incnsv V)]$, and
\code{UpdateWeak} the flux of the weak form
(Section~\ref{sec:incns-weak}).

\textbf{Selection.} Only pressure Neumann regions tagged
\code{USERDEFINEDTYPE="H"} receive the computed flux (\code{IsHOPressure}),
as in the legacy solver; the default Neumann operator of the pressure system
leaves them out. An untagged pressure Neumann region keeps its session
value.

\textbf{Data flow.} Setup (\code{SetVelocityContext}) matches every
selected boundary face to its owning velocity element by geometry id and
builds a restricted expansion list of the boundary-adjacent elements. Each
update then:
\begin{enumerate}
  \item gathers $\incnsv V$ onto the restricted elements
  (\code{GatherBoundaryElements}, an indexed copy) and applies
  \code{CurlCurlOp} there, since the derivatives need the volume polynomial;
  \item forms $\incnsv P-\nu\incnsv C$ on the restricted volume
  (\code{HOPressureVolumeCombineKernel}), using the linearity of the trace;
  \item traces it to the boundary quadrature points with one
  \code{PhysTraceExtractBndCond} per region and contracts with the outward
  normal (\code{HOPressureBoundaryCombineKernel}); on a region with a
  velocity component that is not Dirichlet it also traces
  $\incnsv P-\incnsv R$ and subtracts its contraction with the normal
  components of those velocity components;
  \item subtracts $s\,\incnsv n\cdot\incnsv u_D$
  (\code{SubtractBoundaryVelocity});
  \item forms the boundary inner product
  $\incnsv b_f=B_f^{\mathsf T}W_fJ_f\incnsv g_f$ with \code{IProductWRTBaseOp}
  and reorders it to the pressure block layout
  (\code{BuildBoundaryCoeffs});
  \item in \code{v\_Apply}, scatter-adds
  $\incnsv r_p\mathrel{+}=\sum_fE_f^{\mathsf T}\incnsv b_f$ with orientation
  signs into the Poisson right-hand side. Coefficients shared by several
  boundary pieces are summed in groups, which avoids atomics.
\end{enumerate}
The redesign Poisson operator is positive, so the weak form is
$(\nabla p,\nabla v)=-(f_p,v)+\langle\partial p/\partial n,v\rangle_{\Gamma_N}$
and the boundary load is added, not assigned.

\subsection{Dong outflow condition (HOutflow)}
\label{sec:incns-houtflow}

\textbf{Equations.} On an open boundary $\Gamma_o$ the condition of Dong,
Karniadakis and Chryssostomidis~\cite{DongKC2014} couples a pressure Dirichlet
value, evaluated before the pressure solve, with a velocity Neumann value,
evaluated after it:
\begin{align}
  p &= \nu\,\incnsv n\cdot(\nabla\incnsv u^*)\incnsv n
       - \bigl(\incnsv E(\incnsv n,\incnsv u^*)-\incnsv f_b\bigr)\cdot\incnsv n,
  \label{eq:incns-hout-p}\\
  \frac{\partial\incnsv u^{n+1}}{\partial n} &= \frac{1}{\nu}\Bigl\{
       \bigl[p^{n+1}-\nu\nabla\cdot\incnsv u^*\bigr]\incnsv n
       + \incnsv E(\incnsv n,\incnsv u^*)-\incnsv f_b\Bigr\},
  \label{eq:incns-hout-u}
\end{align}
with the backflow energy flux and smooth switch
\begin{equation}
  \begin{gathered}
  \incnsv E=\Bigl[\frac{\theta+\alpha_2}{2}|\incnsv u^*|^2\incnsv n
  + \frac{1-\theta+\alpha_1}{2}u_n\incnsv u^*\Bigr]S_0(u_n),\\
  S_0(u_n)=\frac12\Bigl[1-\tanh\frac{u_n}{U_0\delta}\Bigr],\qquad
  u_n=\incnsv n\cdot\incnsv u^*.
  \end{gathered}
  \label{eq:incns-hout-e}
\end{equation}
$S_0\approx1$ where the flow locally re-enters ($u_n<0$) and $S_0\approx0$
where it leaves. The vector $\incnsv f_b$ is the velocity \code{VALUE} of the
region. The divergence term in~\eqref{eq:incns-hout-u} is required.

\textbf{Scaling.} \code{UpdatePressure} receives the extrapolated
velocity $\incnsv V^*$ of~\eqref{eq:incns-ustar}, with
$s=1/(\lambda\incnsdt)$ set by \code{SetScale}, and applies $s$ to both the
traced velocity and its gradient, $\incnsv u^*=s\incnsv V^*$ and
$\nabla\incnsv u^*=s\nabla\incnsv V^*$, so that $S_0$, $\incnsv E$, the
normal stress and the divergence use the same field.

\textbf{Stage workflow.}
\begin{enumerate}
  \item \code{UpdatePressure} (before the Poisson solve): gather the
  outlet-adjacent velocity elements, apply \code{PhysDerivOp}, trace
  $\incnsv u^*$ and $\nabla\incnsv u^*$ with \code{PhysTraceExtractBndCond},
  and evaluate~\eqref{eq:incns-hout-p} pointwise
  (\code{HOutflowPressureKernel}). The kernel keeps
  $\incnsv G=\incnsv E-\incnsv f_b$ and $\nabla\cdot\incnsv u^*$ on the
  active backend. The pressure values are projected with \code{FwdTransBCOp}
  and \emph{set} as Dirichlet coefficients by the pressure Dirichlet adapter
  (\code{GetPressureDirichletOp}), after the default Dirichlet operator has
  evaluated the placeholder session value.
  \item After the Poisson solve, \code{SetPressure} hands the operator
  $p^{n+1}$. When the velocity solve applies the velocity Neumann adapter
  (\code{GetVelocityNeumannOp}), the adapter traces $p^{n+1}$ to $\Gamma_o$,
  forms $\incnsv h=[(p-\nu\nabla\cdot\incnsv u^*)\incnsv n+\incnsv G]/\nu$
  (\code{HOutflowVelocityKernel}), projects it with \code{IProductWRTBaseOp}
  and scatter-adds it. The adapter is added with
  \code{AddNeumannBCOp(op, "HOutflow")}, so the default Neumann operator
  leaves the outlet out and $\incnsv h$ is the whole load; $\incnsv f_b$ enters
  through $\incnsv G$.
\end{enumerate}

\textbf{Session.} Every velocity component must carry a Neumann condition
and the pressure a Dirichlet condition, all tagged
\code{USERDEFINEDTYPE="HOutflow"}:

\begin{lstlisting}[style=XMLStyle]
<REGION REF="1">
  <N VAR="u" USERDEFINEDTYPE="HOutflow" VALUE="0" />
  <N VAR="v" USERDEFINEDTYPE="HOutflow" VALUE="0" />
  <D VAR="p" USERDEFINEDTYPE="HOutflow" VALUE="0" />
</REGION>
\end{lstlisting}

\begin{table}[htbp]
\centering
\small
\begin{tabular}{@{}llll@{}}
\toprule
parameter & symbol & default & role \\
\midrule
\code{OutflowBC\_theta}   & $\theta$   & $1$      & weight of the normal kinetic-energy term \\
\code{OutflowBC\_alpha1}  & $\alpha_1$ & $0$      & weight of the $u_n\incnsv u$ term \\
\code{OutflowBC\_alpha2}  & $\alpha_2$ & $0$      & added to the normal kinetic-energy term \\
\code{U0\_HighOrderBC}    & $U_0$      & $1$      & reference velocity of the switch \\
\code{Delta\_HighOrderBC} & $\delta$   & $1/20$   & width of the switch \\
\bottomrule
\end{tabular}
\caption{HOutflow session parameters (read in the
\code{HOutflowBndCondOpImpl} constructor). $U_0$ and $\delta$ must be
positive.}
\label{tab:incns-houtflow}
\end{table}

The pairing of tags, the positivity of $U_0$ and $\delta$, and the
rejection of a tagged pressure Neumann condition are checked with
\code{ASSERTL0}, so also in release builds.

%%%%%%%%%%%%%%%%%%%%%%%%%%%%%%%
\section{Initial conditions and restarts}
\label{sec:incns-ic}

\code{VelocityCorrectionScheme::v\_SetInitialConditions} calls
\code{UnsteadySystem::v\_SetInitialConditions} for the velocity and scalars
and then treats the pressure in the same way, on its own expansion list.
Both go through \code{UnsteadySystem::SetContinuousInitialConditions}, which
brings either source of an initial condition into the continuous space:
\begin{enumerate}
  \item element coefficients: an expression is sampled at the quadrature
  points and projected element by element; a file (\code{.fld}/\code{.chk},
  \code{SessionFunction::EvaluateFld}) is used as read;
  \item every local copy of a global degree of freedom is replaced by the
  signed average of all copies (\code{AvgAssmbScatrOp});
  \item the method selected by \code{SOLVERINFO ProjectInitialConditions}:
  \code{Average} (the default) writes the Dirichlet values at $t_0$ over the
  boundary degrees of freedom (\code{DirBndCondOp}); \code{Galerkin}
  projects the field at the quadrature points with the global mass matrix
  solve of \code{FwdTransOp}, with the Dirichlet values at $t_0$, starting
  from the averaged coefficients.
\end{enumerate}
The start time of a file IC is taken from the file
(\code{CheckForRestartTime}). As in the legacy solver
(\code{AvgAssemble}, \code{GlobalToLocal} and the Dirichlet values imposed
by \code{VelocityCorrectionScheme::v\_DoInitialise}), a file IC that
violates the velocity boundary conditions, for instance one interpolated
from a lower order, therefore starts from velocity Dirichlet values that
match the session; otherwise the first step would correct them and produce
a pressure of order $|\incnsv u-\incnsv u_D|/\incnsdt$
(Section~\ref{sec:incns-kick}).

\textbf{Restarts.} The multistep history is not stored in checkpoints, so a
restart is a cold start from the stored state and repeats the IMEXdirk
start-up step of Section~\ref{sec:incns-time}.

%%%%%%%%%%%%%%%%%%%%%%%%%%%%%%%
\section{Pressure kick at start-up and restart}
\label{sec:incns-kick}

With IMEX order 2, the pressure and the pressure forces ring for a few steps
after every start and restart. On a 2D NACA0012 case ($Re=50\,000$,
$\alpha=5^\circ$, $P=6$, $\incnsdt=2\times10^{-6}$, semi-implicit, strong
form) the ringing was, before the remedies below, about three times larger
than in the legacy solver, which starts IMEX order $J$ with an IMEX order~1
step ($\lambda=1$) instead of IMEXdirk. Two
mechanisms were identified; execution space, operator implementation,
solver tolerances, dealiasing and restart I/O were ruled out.

\textbf{Mechanism 1: weight of the divergence term.} The velocity from
the Helmholtz solve is not discretely divergence-free (on the NACA mesh
$\|\nabla\cdot\incnsv u^n\|_{L^2}\approx0.72$ in every step, in both codes).
Since $\hat{\incnsv u}=\incnsv u^n+O(\incnsdt)$, the pressure contains a
divergence-cleaning part
\begin{equation}
  p_{\mathrm{div}}\approx\frac{1}{\lambda\incnsdt}\,
  \nabla^{-2}\bigl(\nabla\cdot\incnsv u^n\bigr),
  \label{eq:incns-pdiv}
\end{equation}
whose weight changes with the step type (Table~\ref{tab:incns-kick-weight}).
The velocity correction removes the divergence whatever $\lambda$ is, but the
pressure, and with it the wall force, jumps at each start-up and relaxes to
its multistep level. The measured Poisson forcing ratio between DIRK stages
and multistep steps is $2.28=(2/3)/0.2929$. This mechanism dominates on
meshes with thin boundary-layer elements.

\begin{table}[htbp]
\centering
\small
\begin{tabular}{@{}lcc@{}}
\toprule
step type & weight $1/(\lambda\incnsdt)$ & change against multistep \\
\midrule
IMEX-2 multistep & $1.5/\incnsdt$ & --- \\
redesign IMEXdirk 2\_3\_2 start-up (both stages) & $3.41/\incnsdt$ & $+128\%$ \\
legacy IMEX-1 start-up & $1/\incnsdt$ & $-33\%$ \\
\bottomrule
\end{tabular}
\caption{Weight of the divergence part of the pressure.}
\label{tab:incns-kick-weight}
\end{table}

\textbf{Mechanism 2: normalisation of the pressure-BC velocity.} By
\eqref{eq:incns-ustar} the curl--curl velocity of HOPBC and the HOutflow
$\incnsv u^*$ are proper extrapolations only if $\sum_ic_i=\lambda$. In
IMEXdirk 2\_3\_2 stage~2 the weights are $[1-\delta,\delta]$ with
$\sum c=1$, while $\lambda=\gamma\approx0.2929$, so without a normalisation
these velocities would be $1/\gamma\approx3.41$ times too large; remedy A
normalises them. The predictor part is consistent. Without the normalisation
this would also affect IMEXdirk used as the main scheme whenever HOPBC or
HOutflow boundaries are present. The linear-implicit advection has the same
root cause: the right-hand side loses the extrapolated advection with weight
$\sum c\,\incnsdt$, while the advection operator puts it back with weight
$\lambda\incnsdt$, so without a correction stage~2 would keep only
$\gamma\approx29\%$ of its advection. On Kovasznay flow with HOPBC (IMEX-2,
$\incnsdt=0.01$, exact IC), normalising by $\sum c$ reduces the pressure
$L^2$ error after step~1 from $0.169$ to $2.1\times10^{-5}$.

\textbf{Key evidence.} NACA0012 restart test: an uninterrupted run is
compared with a run restarted half-way. The steady pressure drag is
$0.0078$; viscous forces are unaffected.

\begin{table}[htbp]
\centering
\small
\begin{tabular}{@{}lcc@{}}
\toprule
variant & $\max|\Delta F_{\mathrm{drag},p}|$ & $\max|\Delta F_{\mathrm{lift},p}|$ \\
\midrule
legacy & $0.035$ & $5.6\times10^{-3}$ \\
redesign without fixes & $0.097$ & $2.0\times10^{-2}$ \\
fix A only & $0.087$ & $1.6\times10^{-2}$ \\
divergence re-weighting (B1) only & $0.009$ & $7.6\times10^{-3}$ \\
fix A + B1 & $5.5\times10^{-5}$ & $9.9\times10^{-4}$ \\
Euler start-up (B4) & $0.035$ & $5.4\times10^{-3}$ \\
\bottomrule
\end{tabular}
\caption{Deviation of the pressure forces after a restart, on the
NACA0012 case above.}
\label{tab:incns-kick-evidence}
\end{table}

\textbf{Remedies.}
\begin{itemize}
  \item \textbf{A. Normalise the pressure-BC velocity by $\sum c$}
  (implemented): the pressure-BC histories hold $\incnsv u$ and $\incnsv N$
  and are combined with the weights scaled to sum to $\lambda\incnsdt$, the
  scale being $\lambda\incnsdt/\sum c$ with $\sum c$ from
  \code{TimeOp::GetExplicitContributionCoefficientSum}
  (eq.~\eqref{eq:incns-ustar}); it is $\incnsdt$ in multistep steps and DIRK
  stage~1. The linear-implicit advection velocity is scaled by
  $\sum c\,\incnsdt/(\lambda\incnsdt)$, the ratio of the two advection
  weights.
  \item \textbf{B1. Fixed divergence weight during start-up} (not
  implemented): in start-up stage $s$ a fraction
  $\beta=1-\lambda_s/\lambda$ of $\nabla\cdot\incnsv u_s$ would be removed from
  the Poisson forcing ($\beta=0.561,0.201,0.092$ for orders 2, 3, 4), so the
  divergence-cleaning pressure keeps its multistep weight. Here $\incnsv u_s$ is
  the solution the stage starts from, not $\incnsv u^n$: re-weighting $\incnsv
  u^n$ in stage~2 would remove again the part that stage~1 already cleaned.
  \item \textbf{B2. Checkpoint the multistep history} (not implemented):
  restarts continue without a start-up step and reproduce an uninterrupted
  run, but cold starts are unchanged and checkpoints grow.
  \begin{itemize}
    \item \textbf{What is written.} Only the $J-1$ previous velocity levels,
    $\incnsdt$, the order and the step count. The \code{TimeOp} explicit
    history and the three VCS histories (advection RHS, pressure-BC RHS,
    pressure-BC velocity) are deterministic functions of those levels and
    are recomputed on restart. HOPBC and HOutflow keep no step-to-step state.
    \item \textbf{Estimate.} About 400--600 lines plus restart regression
    tests, roughly 3--5 days; moderate-to-low risk. The work splits into a
    generic \code{TimeOp} history export/import (reusing the existing
    \code{SetSolutions}/\code{SetExplicits}/\code{SetStep}), a VCS hook to
    rebuild its histories, checkpoint output of extra fields
    (\code{u\_nm1}, \dots) with metadata, and a restart reader with a
    fallback to the start-up.
    \item \textbf{Constraints.} Exact restarts are bit-exact on CPU only
    (Device reductions are not deterministic). Checkpoints grow by
    $1+(J-1)d/(d+1)$, e.g.\ $+67\%$ in 2D at order 2.
  \end{itemize}
  \item \textbf{B3. Reduce $\nabla\cdot\incnsv u$} (weak pressure form,
  Taylor--Hood spaces, grad--div penalty, extra projection): addresses the
  cause in every step, at the price of larger changes to the method.
  \item \textbf{B4. Legacy-style IMEX-1 start-up} (not implemented):
  reproduces the legacy force history to 3--4 digits, but only reduces the
  kick to the legacy level.
\end{itemize}

\textbf{Steady-state dependence on $\incnsdt$.} A study on Kovasznay flow and
steady cylinder flow ($Re=20,40$) found:
\begin{itemize}
  \item \textbf{Common to both codes.} With HOPBC boundaries, the steady
  pressure error grows mildly as $\lambda\incnsdt$ decreases. The effect
  scales with the spatial error, vanishes spectrally with $P$, and is absent
  with Dirichlet pressure everywhere. $\nabla\cdot\incnsv u$ is a spatial
  error and is not $O(\incnsdt)$.
  \item \textbf{Redesign only (fixed).} An additional $1/(\lambda\incnsdt)$
  error came from subtracting the analytic Dirichlet velocity instead of the
  projected one in the HOPBC flux. The HOPBC now subtracts the discrete
  boundary velocity (as legacy \code{AddVelBC} does), and the redesign matches
  legacy to 3--4 digits.
\end{itemize}
B1, B2 and B4 only affect transients, not the fixed point.

%%%%%%%%%%%%%%%%%%%%%%%%%%%%%%%
\section{Usage and parameters}
\label{sec:incns-usage}

\begin{table}[htbp]
\centering
\small
\begin{tabularx}{\textwidth}{@{}>{\raggedright\arraybackslash}p{0.34\textwidth}
                              >{\raggedright\arraybackslash}X@{}}
\toprule
entry & meaning \\
\midrule
\multicolumn{2}{@{}l}{\textit{SOLVERINFO}} \\
\code{SolverType} & \code{VelocityCorrectionScheme}; selects the equation
system (overrides \code{EQTYPE}). \\
\code{EQTYPE} & required by the driver, e.g.\ \code{UnsteadyNavierStokes};
otherwise only reported in the summary. \\
\code{Projection} & must be \code{Galerkin}. \\
\code{FORMULATION} & \code{SemiImplicit} or \code{LinearImplicit}, case
insensitive. Required: there is no default. \\
\code{SPECTRALHPDEALIASING} & \code{True} (default) or \code{False}. \\
\code{AdvectionForm} & not read; the convective form is always used. \\
\code{PressureFormulation} & \code{Strong} or \code{Weak}
(Section~\ref{sec:incns-weak}), case insensitive. Default \code{Strong}
with \code{SemiImplicit}, \code{Weak} with \code{LinearImplicit}. \\
\code{ProjectInitialConditions} & \code{Average} (default) or
\code{Galerkin} (Section~\ref{sec:incns-ic}). \\
\multicolumn{2}{@{}l}{\textit{TIMEINTEGRATIONSCHEME}} \\
\code{METHOD}, \code{ORDER} & \code{IMEX} with order 1--4. \code{IMEXdirk}
(with \code{VARIANT}) is also available. \\
\multicolumn{2}{@{}l}{\textit{PARAMETERS}} \\
\code{Kinvis} & $\nu$, default 1. \\
\code{TimeStep}, \code{NumSteps}, \code{FinTime} & time step (must be
positive), number of steps, final time. \\
\code{IO\_InfoSteps}, \code{IO\_CFLSteps}, \code{CFL} & cadence of the
step and Courant-number reports, CFL safety factor. \\
\code{D00}, \code{D01}, \dots & anisotropic diffusion tensor, default
identity. \\
HOutflow parameters & Table~\ref{tab:incns-houtflow}. \\
\multicolumn{2}{@{}l}{\textit{GLOBALSYSSOLNINFO (per variable) or PARAMETERS}} \\
\code{IterativeSolverTolerance} & default $10^{-9}$. \\
\code{NekLinSysMaxIterations} & default 5000. \\
\code{AbsoluteTolerance} & \code{True} for an absolute criterion
(GLOBALSYSSOLNINFO only). \\
\code{Preconditioner} & \code{Diagonal} (default) or \code{Null}; also
SOLVERINFO. \\
\code{LinSysIterSolver} & ignored: the pressure, velocity and
initial-condition projection solves use \code{ConjGrad}, or \code{GMRES} for
the linear-implicit velocity. \\
\multicolumn{2}{@{}l}{\textit{command line}} \\
\code{--opExecSpace} & \code{Serial}, \code{AVX} or \code{Device}; default
\code{Device} if a device is present, else \code{AVX} or \code{Serial}, with a
warning. \\
\code{--opImpl} & \code{StdMat}, \code{SumFac} (default) or
\code{SumFacTOP} (\code{Device} only); overrides per-operator choices in
\code{OPTIMISATION/OPERATORS}. \\
\code{-I}, \code{-P} & override a SOLVERINFO entry or a parameter, e.g.\
\code{-I Formulation=semiimplicit}. \\
\code{-v} & verbose output, including linear-solver iteration counts. \\
\bottomrule
\end{tabularx}
\caption{Session and command-line settings used by the solver.}
\label{tab:incns-params}
\end{table}

\textbf{Linear-solver settings.} Each solve reads the GLOBALSYSSOLNINFO
entry of its first variable, falling back to the parameter of the same name:
the velocity components are solved together and use the setting of
\code{u}, the pressure that of \code{p}. With the default relative criterion
an iteration stops when $\|\incnsv r\|<\epsilon\,\|\incnsv b-A\incnsv x_0\|$,
with the previous solution as initial guess $\incnsv x_0$. With
\code{AbsoluteTolerance=True} it stops when $\|\incnsv r\|<\epsilon$. The
right-hand sides are scaled: the velocity forcing~\eqref{eq:incns-helm}
contains $\hat{\incnsv u}/(\nu\lambda\incnsdt)$ and the pressure forcing
$\nabla\cdot\hat{\incnsv u}/(\lambda\incnsdt)$. For small $\nu$ and
$\incnsdt$ these are many orders of magnitude larger than the velocity, so
an absolute tolerance on them is far stricter than it looks and can prevent
convergence.

\textbf{Example.} A minimal conditions block, adapted from
\path{Tests/KovasznayFlow2D.xml}: the test sets \code{ORDER} to 1 and passes
the formulation on the command line (\code{-I Formulation=...}).

\begin{lstlisting}[style=XMLStyle]
<SOLVERINFO>
  <I PROPERTY="SolverType"  VALUE="VelocityCorrectionScheme" />
  <I PROPERTY="EQTYPE"      VALUE="UnsteadyNavierStokes" />
  <I PROPERTY="Projection"  VALUE="Galerkin" />
  <I PROPERTY="Formulation" VALUE="SemiImplicit" />
</SOLVERINFO>
<TIMEINTEGRATIONSCHEME>
  <METHOD> IMEX </METHOD>
  <ORDER> 2 </ORDER>
</TIMEINTEGRATIONSCHEME>
<PARAMETERS>
  <P> TimeStep = 0.01 </P>
  <P> NumSteps = 10 </P>
  <P> IO_InfoSteps = NumSteps </P>
  <P> Kinvis = 0.025 </P>
  <P> IterativeSolverTolerance = 1e-12 </P>
</PARAMETERS>
\end{lstlisting}

A wall with a higher-order pressure condition is declared as
\lstinline[style=XMLStyle]!<N VAR="p" USERDEFINEDTYPE="H" VALUE="0" />!
next to Dirichlet velocity conditions (see
\path{Tests/ChannelFlowHOPBC2D.xml}). The solver is run as
\begin{lstlisting}[style=BashInputStyle]
IncNavierStokesSolverRedesign --opExecSpace=Serial --opImpl=SumFac \
    KovasznayFlow2D.xml
\end{lstlisting}

%%%%%%%%%%%%%%%%%%%%%%%%%%%%%%%
\section{Implementation map}
\label{sec:incns-map}

\begin{table}[htbp]
\centering
\small
\begin{tabularx}{\textwidth}{@{}>{\raggedright\arraybackslash}p{0.47\textwidth}
                              >{\raggedright\arraybackslash}X@{}}
\toprule
location & role \\
\midrule
\multicolumn{2}{@{}l}{\path{solvers/IncNavierStokesSolverRedesign/}} \\
\path{EquationSystems/VelocityCorrectionScheme.cpp}:
\code{SolveUnsteadyStokesSystem} & pressure forcing, BC updates, Poisson and
velocity solves (Section~\ref{sec:incns-vcs}) \\
\quad\code{EvaluateAdvection\_SetPressureBCs} & explicit advection and
forcing, refresh of the extrapolation histories \\
\quad\code{SetUpWeakPressureForcing} & volume term of the weak pressure
form (Section~\ref{sec:incns-weak}) \\
\quad\code{InitialiseParameters}, \code{v\_InitialiseOperators} &
\code{FORMULATION}, \code{PressureFormulation}, dealiasing; operator and
solver set-up, BC registration \\
\quad\code{v\_SetInitialConditions} & velocity and pressure ICs
(Section~\ref{sec:incns-ic}) \\
\path{BoundaryConditions/HOPressureBndCond/HOPressureBndCondOpImpl.hpp} &
HOPBC: \code{BuildBoundaryData} (tag selection),
\code{BuildBoundaryElementData} (restricted list, extract operators),
\code{BuildBoundaryVelocityData} (normals, discrete Dirichlet velocity),
\code{GatherBoundaryElements}, \code{BuildBoundaryPhys},
\code{EvaluateBoundaryVelocity}, \code{SubtractBoundaryVelocity},
\code{BuildBoundaryCoeffs}, \code{v\_Apply} \\
\path{BoundaryConditions/HOutflowBndCond/HOutflowBndCondOpImpl.hpp} &
HOutflow: \code{HOutflowPressureKernel}, \code{HOutflowVelocityKernel},
\code{ValidateAndBuildVelocityBoundaryData}, \code{EvaluatePressure},
\code{EvaluateVelocity}, \code{ApplyNeumannMaps}, and the adapters
\code{HOutflowPressureAdapter} and \code{HOutflowVelocityAdapter} \\
\code{ADD\_OPERATOR} in \path{CMakeLists.txt} & per execution space
factory registration \\
\path{UnitTests/TestHOPressureBndCond.cpp},
\path{UnitTests/TestHOutflowBndCond.cpp} & operator unit tests \\
\path{Tests/} & regression sessions, e.g.\ \code{KovasznayFlowHOPBC2D},
\code{KovasznayFlowHOPBC2D\_SmallTimeStep},
\code{KovasznayFlowHOPBC2D\_IMEXOrder2}, \code{KovasznayFlowHOPBC2D\_IMEXdirk},
\code{HOutflowKovasznay2D}, \code{HOutflowBackflow2D}, and the
\code{\_WeakPressure} variants \\
\multicolumn{2}{@{}l}{\path{library/SolverCore/}} \\
\path{TimeOps/IMEX/IMEXOpImpl.hpp} & IMEX multistep scheme and its IMEXdirk
start-up \\
\path{TimeOps/IMEXdirk/IMEXdirkOpImpl.hpp} & DIRK stages and their
explicit-contribution coefficients \\
\path{TimeOps/TimeOp.hpp} & \code{ExtrapolateExplicitContribution},
\code{LatestExplicitContribution},
\code{GetExplicitContributionCoefficientSum} \\
\path{EquationSystems/UnsteadySystem.cpp} & time-stepping parameters,
filters, \code{SetContinuousInitialConditions}, restart time \\
\path{GlobalLinSysOps/LinearSystems/PoissonSolve},
\path{HelmSolve}, \path{LinearADRSolve} & pressure and velocity systems;
\code{PoissonSolveOp::ApplyRhs} takes a coefficient-space right-hand side \\
\path{GlobalLinSysOps/LinearSolvers/LinearSolverOp.hpp} & solver selection,
tolerances \\
\path{MeanRemoval/MeanRemovalOp.hpp} & zero-mean forcing of a singular
strong pressure system, and of the reported pressure error \\
\multicolumn{2}{@{}l}{\path{library/MultiRegions/}} \\
\path{ElmtOps/CurlCurl}, \path{ElmtOps/IProductWRTDerivBase},
\path{ElmtOps/PhysTraceExtract}, \path{BndCondOps/} & curl--curl, weak
pressure volume term, volume-to-boundary extraction, Dirichlet,
Neumann, Robin and \code{FwdTransBC} operators \\
\bottomrule
\end{tabularx}
\caption{Where the solver's pieces live.}
\label{tab:incns-map}
\end{table}

\fussy
